\chapter{Delta-One Instruments}\label{ch:m1:delta-one-instruments}

Five investors want the same thing: the return of one equity index on a
hundred million for a year. The first buys the shares, the second a fund, the
third futures, the fourth signs a swap with a bank, the fifth opens a
\omterm{def:m1:delta-one-instruments:cfd}{contract for difference} with a broker. At the end of the year the index has
returned what it returned, and the five have earned five different amounts.
The most expensive route cost more than five times the cheapest, and which route is
cheapest depends on who the investor is: whether it has the cash, what tax
it pays on dividends, how long it holds. None of the five instruments
contains an option, a view or a model. They differ in financing, in tax and in
plumbing, which is to say in everything the first sixteen chapters were
about.

\section{What delta-one means}

\begin{definition}[Delta-one]\label{def:m1:delta-one-instruments:deltaone}
An instrument is \emph{delta-one}\index{delta-one} on an underlying when its
value changes by one unit for a one-unit change in the value of the
underlying: it carries the exposure without optionality. Shares, funds,
futures and forwards, \omterm{def:m1:financing:trs}{total return swaps}, contracts for difference,
\omterm{def:m1:delta-one-instruments:dr}{depositary receipts} and certificates that track an asset one for one are
delta-one. A \emph{delta-one desk} prices, hedges and arbitrages the
differences between them.
\end{definition}

Since all wrappers deliver the same exposure, only their costs can differ.
There are five kinds of cost, and every wrapper is a vector of five numbers.

\begin{method}[The five costs of a wrapper]\label{met:m1:delta-one-instruments:five}
\begin{enumerate}
\item \textbf{Trading}: commissions, spread and impact to enter and leave,
and to roll for instruments that expire.
\item \textbf{Transaction tax} on the purchase, where one exists.
\item \textbf{Running fees}: a fund's \omterm{def:m1:the-buy-side:fees}{management fee}.
\item \textbf{Funding}: the spread over the benchmark rate paid on the money
that is borrowed, explicitly (a margin loan) or inside the price (a future,
a swap).
\item \textbf{Dividends}: the fraction lost to \omterm{def:m1:delta-one-instruments:wht}{withholding tax}, which depends
on the holder \emph{and} on the wrapper.
\end{enumerate}
Express each in basis points of notional over the intended holding period,
and add.
\end{method}

\section{Swaps, contracts for difference and the funding spread}

\begin{definition}[Equity swap and funding spread]\label{def:m1:delta-one-instruments:swap}
An \emph{equity swap}\index{equity swap} is a \omterm{def:m1:financing:trs}{total return swap}
(\cref{def:m1:financing:trs}) on a share, a basket or an index, with periodic
resets at which the accrued performance is paid and the notional is reset to
the current price. Its \emph{funding spread}\index{funding spread} is the
margin over the benchmark rate in the financing leg: the price of the
dealer's balance sheet, and the number on which swaps compete.
\end{definition}

\begin{omfigure}
\begin{tikzpicture}[font=\small,
  box/.style={draw=omDef,rounded corners=2pt,minimum width=2.2cm,minimum height=1.1cm,align=center,fill=omDef!4},
  flow/.style={-{Stealth[length=2mm]},thick}]
\node[box] (c) at (0,0) {Client};
\node[box] (d) at (7.0,0) {Dealer};
\node[box,draw=omExo,fill=omExo!6] (m) at (12.2,0) {Market};
\draw[flow,omDef] ([yshift=9pt]d.west) -- node[above,font=\footnotesize,align=center]{price return\\$+$ a share $\eta$ of the dividends} ([yshift=9pt]c.east);
\draw[flow,omThm] ([yshift=-9pt]c.east) -- node[below,font=\footnotesize,align=center]{benchmark rate\\$+$ funding spread} ([yshift=-9pt]d.west);
\draw[flow,omExo] (d.east) -- node[above,font=\footnotesize]{buys the shares} node[below,font=\footnotesize]{as its hedge} (m.west);
\end{tikzpicture}

\omcaption{An \omterm{def:m1:delta-one-instruments:swap}{equity swap}. The dealer is flat: it owns what it owes. Its
profit is the \omterm{def:m1:delta-one-instruments:swap}{funding spread} less its own cost of financing the shares, plus
whatever the shares earn in its hands (lending fees, a better dividend tax
position) that it does not pass on.}\label{fig:m1:delta-one-instruments:swap}
\end{omfigure}

\begin{definition}[Contract for difference]\label{def:m1:delta-one-instruments:cfd}
A \emph{contract for difference}\index{contract for difference} (CFD) is an
open-ended \omterm{def:m1:delta-one-instruments:swap}{equity swap} offered by a broker to its clients, usually retail,
on margin: the client receives the change in price, is credited or debited
dividends, and pays daily financing on the full notional for as long as the
position is open.
\end{definition}

A swap and a CFD are the same contract at two price points. The institution
negotiates a \omterm{def:m1:delta-one-instruments:swap}{funding spread} of tens of basis points with several dealers; the
retail client takes the broker's rate, commonly a few \emph{percent} over
the benchmark. Both exist for the same two reasons: leverage without a margin
loan, and the fact that the client never owns the share.

\begin{definition}[Withholding tax and dividend enhancement]\label{def:m1:delta-one-instruments:wht}
A \emph{withholding tax}\index{withholding tax} is a tax deducted at source
from dividends paid to a non-resident holder, at a statutory rate that a tax
treaty may reduce. \emph{Dividend enhancement}\index{dividend enhancement} is
any arrangement in which the exposure is held, over the dividend date, by the
party with the lowest tax cost, who passes part of the saving to the
economic holder, typically through the share $\eta$ of the dividend paid in
a swap.
\end{definition}

\begin{dated}{2026-09}{Two taxes that shape wrappers}\label{dat:m1:delta-one-instruments:tax}
\textbf{United Kingdom.} Buying UK shares electronically costs 0.5\% in Stamp
Duty Reserve Tax; transferring shares into some depositary-receipt schemes or
clearance services costs 1.5\%; foreign shares bought outside the UK are
normally not charged. Recognised intermediaries such as \omterm{def:m1:what-a-trading-firm-does:market-maker}{market makers} are
relieved of the tax on their purchases, which is what lets a dealer hedge a
swap or a CFD, on which no share changes hands, without paying it.

\textbf{United States.} Section 871(m) of the tax code treats
\emph{dividend-equivalent} payments on derivatives over US shares as US-source
dividends, subject to the same withholding as the dividends themselves. It
applies to \omterm{def:m1:delta-one-instruments:deltaone}{delta-one} contracts; under Notice 2024-44 contracts that are not
\omterm{def:m1:delta-one-instruments:deltaone}{delta-one} and are issued before 1 January 2027 remain outside it. For a
non-US holder, a swap on a US stock no longer improves on the stock.
\end{dated}

\section{Futures as delta-one}

A future (Part III) is the most standardised wrapper: no counterparty to
negotiate with, margin of a few percent, and the financing inside the price.
If the index is at $S$, the benchmark rate is $r$ and the dividend yield $q$,
the future for date $T$ is fair at $S(1+(r-q)T)$; the rate that makes the
\emph{market} price fair is the future's \emph{implied financing rate}, and
its excess over the benchmark is the \omterm{def:m1:delta-one-instruments:swap}{funding spread} of the futures market.
It is set by supply and demand for leverage: when many investors want to be
long through futures and dealers' balance sheets are scarce, as at
year-ends, the spread rises. The holder also pays a roll, four times a year,
whose cost is that spread for the next quarter plus the trading cost of the
calendar spread (\cref{ch:m1:basis-roll-and-delivery}).

\begin{example}[Reading the spread from a price]\label{ex:m1:delta-one-instruments:implied}
Index 6\,000, $r = 4.2\%$, $q = 1.3\%$, three months: fair value $6\,000
\times (1 + 0.029 \times 0.25) = 6\,043.5$. The future trades at 6\,048. Its
implied financing rate is $(6\,048/6\,000 - 1)/0.25 + 1.3\% = 4.5\%$: a
spread of 30 basis points. An investor with cash who buys the future and
keeps the cash at the benchmark pays 30 basis points a year for the
exposure; an investor who would otherwise borrow from a \omterm{def:m1:financing:pb}{prime broker} at the
benchmark plus 50 \emph{saves} 20.
\end{example}

\section{Depositary receipts and dual listings}

\begin{definition}[Depositary receipt]\label{def:m1:delta-one-instruments:dr}
A \emph{depositary receipt}\index{depositary receipt} is a security issued by
a depositary bank that represents a fixed number (the \emph{ratio}) of shares
of a foreign company held in custody in the home market. It trades and
settles in the host market's currency and system; American depositary
receipts (ADRs) are the US form. Receipts can be \emph{issued} against a
deposit of ordinary shares and \emph{cancelled} for delivery of the shares,
for a fee per receipt.
\end{definition}

\begin{definition}[Dual listing]\label{def:m1:delta-one-instruments:dual}
A company has a \emph{dual listing}\index{dual listing} when the same class
of share is listed on exchanges in two countries and can be moved between
the two registers. Unlike a receipt there is no intermediate security, only
two places where the same share trades.
\end{definition}

\begin{dated}{2012-08}{What the regulator's bulletin says about ADR costs}\label{dat:m1:delta-one-instruments:adr}
Depositary banks charge holders a custody fee (depositary services fee),
commonly subtracted from the gross dividend; fees are typically assessed per
ADR, the bulletin's example being \$20 to \$50 for 1\,000 ADRs. The fee
schedule of each programme is in its Form F-6. Sponsored programmes have
three levels: Level 1 trades over the counter, Level 2 is exchange-listed,
Level 3 can raise capital. Unsponsored ADRs are set up without the company's
cooperation.
\end{dated}

\begin{proposition}[The conversion band of a receipt]\label{prop:m1:delta-one-instruments:band}
Let $P^* = \rho\,S\,X$ be the parity price of a receipt with ratio $\rho$,
ordinary price $S$ and exchange rate $X$. With trading costs $c$ (both legs
and the currency trade, as a fraction), issue and cancellation fees $f_i$ and
$f_c$ per receipt, and a tax $\tau$ on depositing shares, conversion is
profitable only when the receipt's premium to parity lies outside
\[
\Bigl[-\bigl(c + f_c/P^*\bigr),\; c + \tau + f_i/P^*\Bigr].
\]
A fee fixed in cents makes the band wider for low-priced receipts, and a
deposit tax makes it asymmetric.
\end{proposition}
\begin{proof}
Above parity: buy $\rho$ ordinaries, pay $\tau$ and $f_i$, receive a receipt,
sell it; the profit per receipt is $P - P^*(1+c+\tau) - f_i$. Below: buy the
receipt, pay $f_c$, receive and sell the ordinaries.
\end{proof}

\begin{omfigure}
\begin{tikzpicture}
\begin{axis}[width=10cm,height=4.8cm,
  xlabel={price of the receipt (\$)},ylabel={premium to parity (bp)},
  tick label style={font=\small},label style={font=\small},
  xmin=0,xmax=105,ymin=-120,ymax=120,grid=major,grid style={gray!25},table/col sep=comma]
\addplot[omDef,very thick,mark=*,mark size=1.5pt] table[x=receipt_price,y=upper_bp]{figdata/markets-1/17-delta-one-instruments/band.csv};
\addplot[omDef,very thick,mark=*,mark size=1.5pt] table[x=receipt_price,y=lower_bp]{figdata/markets-1/17-delta-one-instruments/band.csv};
\node[font=\footnotesize,anchor=west] (in) at (axis cs:62,-72) {inside: nothing pays};
\draw[-{Stealth[length=1.6mm]}] (in.north) -- (axis cs:80,-2);
\node[font=\footnotesize,anchor=west] at (axis cs:32,72) {issue receipts};
\node[font=\footnotesize,anchor=west] at (axis cs:22,-72) {cancel receipts};
\end{axis}
\end{tikzpicture}

\omcaption{The conversion band for fees of 5 cents per receipt each way and
trading costs of 6 basis points, with no deposit tax. At \$5 the fee alone is
100 basis points. Data:
\cref{prop:m1:delta-one-instruments:band}.}\label{fig:m1:delta-one-instruments:band}
\end{omfigure}

Parity is only observable while both markets are open. When the home market
is closed the receipt is the only live price, and its ``premium'' to the
stale home close is information, as for the ETFs of
\cref{prop:m1:exchange-traded-funds:stale}. Some pairs never converge at all:
where conversion is restricted by capital controls or ownership limits, two
lines of the same company can trade tens of percent apart for years, and the
spread is a position, not an arbitrage.

\section{Comparing wrappers}

\Cref{fig:m1:delta-one-instruments:stack} applies
\cref{met:m1:delta-one-instruments:five} to the chapter's five investors for
one year, with the illustrative terms of the tutorial: a market with a 50
basis point purchase tax, a dividend yield of 2\% of which a foreign holder
loses 15\%, a prime-broker spread of 50 basis points, a futures spread of 30, a
swap at 40 with 95\% of the dividend, a fund charging 7, and a retail CFD at
250.

\begin{omfigure}
\begin{tikzpicture}
\begin{axis}[width=11.5cm,height=5.6cm,ybar stacked,bar width=13pt,
  ylabel={cost over one year (bp)},
  symbolic x coords={shares,ETF,future,swap,CFD},xtick=data,
  tick label style={font=\small},label style={font=\small},xticklabel style={font=\small\strut},
  ymin=0,ymax=330,ymajorgrids,grid style={gray!25},enlarge x limits=0.12,
  legend columns=5,legend style={font=\footnotesize,at={(0.5,-0.18)},anchor=north,draw=none,fill=none,/tikz/every even column/.append style={column sep=6pt}},
  table/col sep=comma]
\addplot[fill=omExo!50,draw=omExo] table[x=wrapper,y=trading]{figdata/markets-1/17-delta-one-instruments/stack_leveraged.csv};
\addplot[fill=omThm!70,draw=omThm] table[x=wrapper,y=tax]{figdata/markets-1/17-delta-one-instruments/stack_leveraged.csv};
\addplot[fill=black!60,draw=black] table[x=wrapper,y=running]{figdata/markets-1/17-delta-one-instruments/stack_leveraged.csv};
\addplot[fill=omDef!60,draw=omDef] table[x=wrapper,y=funding]{figdata/markets-1/17-delta-one-instruments/stack_leveraged.csv};
\addplot[fill=omDef!15,draw=omDef] table[x=wrapper,y=dividends]{figdata/markets-1/17-delta-one-instruments/stack_leveraged.csv};
\legend{trading,tax,fee,funding,dividends}
\end{axis}
\end{tikzpicture}

\omcaption{One year of the same exposure for a leveraged foreign investor:
140, 93, 68, 56 and 296 basis points. An investor with cash pays no funding
on shares or the fund, and the order changes: fund 43, swap 56, future 68,
shares 90, CFD 296. Data: the tutorial's comparator, illustrative
terms.}\label{fig:m1:delta-one-instruments:stack}
\end{omfigure}

\begin{omfigure}
\begin{tikzpicture}
\begin{axis}[width=11cm,height=5.2cm,
  xlabel={holding period (years)},ylabel={cost per year (bp)},
  tick label style={font=\small},label style={font=\small},
  xmin=0,xmax=5,ymin=0,ymax=200,grid=major,grid style={gray!25},
  legend columns=4,legend style={font=\footnotesize,at={(0.5,-0.3)},anchor=north,draw=none,fill=none,/tikz/every even column/.append style={column sep=8pt}},
  table/col sep=comma]
\addplot[omThm,very thick] table[x=years,y=shares]{figdata/markets-1/17-delta-one-instruments/horizon.csv};
\addplot[black,thick,dashdotted] table[x=years,y=ETF]{figdata/markets-1/17-delta-one-instruments/horizon.csv};
\addplot[omDef,very thick,dashed] table[x=years,y=future]{figdata/markets-1/17-delta-one-instruments/horizon.csv};
\addplot[omExo,very thick,dotted] table[x=years,y=swap]{figdata/markets-1/17-delta-one-instruments/horizon.csv};
\legend{shares,ETF,future,swap}
\end{axis}
\end{tikzpicture}

\omcaption{Annualised cost against holding period for an investor with cash.
Up-front costs (the purchase tax on shares) are amortised; running costs
(funding inside futures and swaps) are not. Shares overtake the swap after
2.7 years and the future after 1.6; the fund, exempt here from the tax, is
cheapest beyond seven weeks. Data: the tutorial's
comparator.}\label{fig:m1:delta-one-instruments:horizon}
\end{omfigure}

The comparison is mechanical once the five numbers are known; the work is
in knowing them. Three are quoted (trading, fees, the \omterm{def:m1:delta-one-instruments:swap}{funding spread}); the
tax cost depends on the investor's domicile and status; and the dividend
share $\eta$ of a swap is the dealer's private estimate of what the shares
earn in its hands. A \omterm{def:m1:delta-one-instruments:deltaone}{delta-one desk} makes its living from differences of a
few basis points in these numbers between clients, and a quantitative
investor loses a comparable amount by never asking.

\section{Tutorial: five wrappers, one exposure}\label{tut:m1:delta-one-instruments:five}

\begin{tutorial}
\textbf{Goal.} Compute the all-in cost of five wrappers for two investors
and any horizon, and the conversion band of a receipt. \textbf{End state:}
the three data figures of this chapter.
\begin{tutsteps}
\item \textbf{Terms.} One record per wrapper, five kinds of cost.
\omcode{code/markets-1/17-delta-one-instruments/python/wrappers.py}{19}{25}{Illustrative terms: entry, exit, purchase tax, running fee, funding spread, roll cost, rolls per year, dividend leakage.}\label{lst:m1:delta-one-instruments:terms}
\item \textbf{The breakdown.} Synthetic wrappers finance the whole notional
inside their price, whoever holds them.
\omcode{code/markets-1/17-delta-one-instruments/python/wrappers.py}{28}{41}{The five costs over a holding period.}\label{lst:m1:delta-one-instruments:breakdown}
\item \textbf{The receipt band.}
\omcode{code/markets-1/17-delta-one-instruments/python/wrappers.py}{72}{78}{Premiums beyond which issuing or cancelling receipts pays.}\label{lst:m1:delta-one-instruments:band}
\end{tutsteps}
\textbf{What to change next.} Add the short side: a short future
\emph{earns} the rich financing, a short fund earns the fee and pays a
borrow. Then make the futures spread seasonal, higher over the year-end
roll, and see for which horizons the future stops being competitive.
\end{tutorial}

\section{Build: the wrapper comparator}\label{bld:m1:delta-one-instruments:comparator}

\begin{build}
\bfield{Purpose} Before the miniature firm puts on any index or
single-stock exposure, long or short, it asks which wrapper is cheapest for
its horizon and tax position, and records why.
\bfield{Interface} \texttt{WrapperTerms(\ldots)} with a \texttt{synthetic}
flag and a \texttt{shortable} flag; \texttt{Investor(financed,
dividend\_leak\_override)}; \texttt{cost(wrapper, investor, years,
div\_yield\_bp, side, borrow\_fee\_bp)} returning a total and its parts;
\texttt{rank(\ldots)} returning the quotes sorted by total.
\bfield{Rules} Long: purchase tax, funding on the financed fraction
(everything, for synthetics), dividend leakage. Short: no purchase tax, the
short-side spread plus the \omterm{def:m1:financing:seclending}{borrow fee} from
\cref{bld:m1:stock-loan-and-short-selling:borrow}, fund fees change sign, and
wrappers that cannot be shorted are excluded, with an error if one is
requested explicitly. Ties are broken by name.
\bfield{Acceptance tests} \texttt{code/firm/wrappers/tests/}: the chapter's
two rankings; a tax-exempt holder over five years; the short side.
\bfield{Stretch} Replace each point estimate by a range and report the
wrappers that are cheapest somewhere in the box: usually two, which is the
honest answer.
\end{build}

\begin{omsources}
\item US Securities and Exchange Commission, \emph{Investor Bulletin:
American Depositary Receipts}, August 2012.
\item HM Revenue \& Customs, \emph{Tax when you buy shares}, gov.uk; Stamp
Taxes on Shares Manual, STSM042050 (intermediary relief).
\item Internal Revenue Service, Notice 2024-44 (section 871(m) phase-in).
\item A.~Madhavan, U.~Marchioni, W.~Li and D.~Y.~Du, ``Equity ETFs versus
index futures: a comparison for fully funded investors'', \emph{Journal of
Index Investing}, 2014.
\end{omsources}

\section{Exercises}

\begin{exercise}[$\star$]\label{exo:m1:delta-one-instruments:1}
A fund receives the total return of an index on \$50 million and pays the
benchmark, 4.2\%, plus 40 basis points, actual/360. Over a 91-day period the
index returns 6\% with dividends. Give the two payments and the net.
\end{exercise}

\begin{exercise}[$\star$]\label{exo:m1:delta-one-instruments:2}
Reproduce \cref{ex:m1:delta-one-instruments:implied} and give the future's
price at which its \omterm{def:m1:delta-one-instruments:swap}{funding spread} would be zero, and the price at which it
would be 50 basis points.
\end{exercise}

\begin{exercise}[$\star$]\label{exo:m1:delta-one-instruments:3}
An ADR represents 2 ordinary shares. The ordinary trades at \euro42.10, the
euro at \$1.17, the ADR at \$98.90. Give the parity price and the premium in
basis points.
\end{exercise}

\begin{exercise}[$\star\star$]\label{exo:m1:delta-one-instruments:4}
With the chapter's terms, an investor with cash compares shares and the CFD.
Below what holding period is the CFD cheaper? (The dividend leakage is the
same for both.)
\end{exercise}

\begin{exercise}[$\star\star$]\label{exo:m1:delta-one-instruments:5}
A US stock yields 1.5\%. Give the annual cost of withholding for a foreign
holder at the statutory 30\% and at a treaty rate of 15\%. A dealer offers a
swap paying 100\% of the dividend. Using
\cref{dat:m1:delta-one-instruments:tax}, explain why the offer should make the
holder suspicious.
\end{exercise}

\begin{exercise}[$\star\star$]\label{exo:m1:delta-one-instruments:6}
For the investor with cash, verify the two crossovers quoted in
\cref{fig:m1:delta-one-instruments:horizon} (shares against the swap, fund
against the future) from the terms in the tutorial.
\end{exercise}

\begin{exercise}[$\star\star\star$]\label{exo:m1:delta-one-instruments:7}
\emph{Coding.} Run the comparator for the leveraged investor over three
months. Report the five totals and their annualised values. Which wrappers
tie, and what would break the tie in practice?
\end{exercise}

\begin{exercise}[$\star\star\star$]\label{exo:m1:delta-one-instruments:8}
\emph{Find the flaw.} A broker's page says: ``CFDs: no commission, no \omterm{def:m1:asian-and-emerging-equity-markets:stamp}{stamp
duty}. The cheapest way to hold UK shares.'' A client holds for three years,
with the chapter's terms. Compute both costs and name what the page left out.
\end{exercise}

\section{Problem: The Receipt and the Ordinary}

\begin{problem}[{Weekend problem --- one company, two tickers, one ocean}]\label{pb:m1:delta-one-instruments:1}
A UK company's ordinary share trades in London at 500p. Its ADR represents 4
ordinaries and trades in New York. Sterling is at \$1.25. The depositary
charges 5 cents per ADR to issue or to cancel. Depositing UK shares into the
ADR programme costs the 1.5\% charge of
\cref{dat:m1:delta-one-instruments:tax}. Trading both legs and the currency
costs 6 basis points in all.

\medskip\noindent\textbf{Part I --- Parity and band.}
\begin{enumerate}
\item Give the parity price of the ADR.
\item Give the upper bound of the conversion band in basis points.
\item Give the lower bound.
\item The ADR trades at \$25.30. Give its premium. Does any conversion pay?
\item Explain in one sentence why the band is asymmetric, and what that
predicts about the sign of the average premium.
\end{enumerate}

\medskip\noindent\textbf{Part II --- Two conversions.}
\begin{enumerate}[resume]
\item The ADR is at \$25.50. A firm issues 100\,000 ADRs. List the cash flows
and give the profit in dollars and in basis points.
\item The ADR is at \$24.90. The firm cancels 100\,000 ADRs. Same questions.
\item Conversion takes two days. The firm leaves the currency unhedged; with
a sterling volatility of 8\% a year, give the standard deviation of the
currency effect over two days, in basis points. Conclude.
\item Which of the two conversions needs a stock borrow if the firm wants to
sell the leg it will receive \emph{before} it receives it?
\end{enumerate}

\medskip\noindent\textbf{Part III --- After London closes.}
London closes at 11:30 New York time. At 15:00 the ADR trades at \$25.45
while the last London price still gives a parity of \$25.00. The US market
has risen 1.5\% since 11:30, and the stock's beta to it is 1.1.
\begin{enumerate}[resume]
\item Give an updated estimate of parity.
\item Give the ADR's premium to the stale parity and to your estimate.
\item A colleague wants to ``sell the 180 basis point premium''. What is he
actually selling, and against what?
\item How would you test, on history, whether the ADR's afternoon move
predicts the next London open?
\item Why is that predictability not free money for a London opening
auction trader?
\end{enumerate}

\medskip\noindent\textbf{Part IV --- Judgement.}
\begin{enumerate}[resume]
\item For a US investor wanting this company for five years, list what the
ADR costs that the ordinary would not, and the reverse.
\item The company moves its primary listing to New York and the ADR is
replaced by the share itself. What happens to the band?
\item Two classes of the same company trade 30\% apart in two countries
between which conversion is forbidden. What kind of trade is ``buy the cheap
one, sell the dear one''?
\item Who are the natural operators of the conversions of Part II?
\item State the \emph{named result}: the conversion band, in basis points of
parity.
\item In one sentence: what does a \omterm{def:m1:delta-one-instruments:deltaone}{delta-one desk} sell?
\end{enumerate}
\end{problem}

\section{Interview questions}

\begin{interviewq}[$\star$ \iqroles{trader, bank, researcher}]\label{iq:m1:delta-one-instruments:1}
What does ``\omterm{def:m1:delta-one-instruments:deltaone}{delta-one}'' mean, and what does a \omterm{def:m1:delta-one-instruments:deltaone}{delta-one desk} do?
\end{interviewq}

\begin{interviewq}[$\star$ \iqroles{trader, bank}]\label{iq:m1:delta-one-instruments:2}
A client can buy an index fund or go long index futures. How do you compare
the two for her?
\end{interviewq}

\begin{interviewq}[$\star\star$ \iqroles{trader, researcher}]\label{iq:m1:delta-one-instruments:3}
Index futures are trading ``rich'': the implied financing rate is 80 basis
points over the benchmark. Who is hurt, who benefits, and what trade
collects the spread?
\end{interviewq}

\begin{interviewq}[$\star\star$ \iqroles{bank, trader}]\label{iq:m1:delta-one-instruments:4}
You are a dealer quoting a \omterm{def:m1:financing:trs}{total return swap} on a European stock to a \omterm{def:m1:what-a-trading-firm-does:hedge-fund}{hedge
fund}. What goes into your \omterm{def:m1:delta-one-instruments:swap}{funding spread} and your dividend percentage?
\end{interviewq}

\begin{interviewq}[$\star\star$ \iqroles{trader, researcher}]\label{iq:m1:delta-one-instruments:5}
An ADR trades 1\% above its parity. Is that an arbitrage? What do you check?
\end{interviewq}

\begin{interviewq}[$\star\star\star$ \iqroles{researcher, trader}]\label{iq:m1:delta-one-instruments:6}
Your equity long-short backtest assumes you hold shares. The firm would
implement it on swap. What changes in the backtest?
\end{interviewq}
